% 讨论与结论
\chapter{讨论与结论}\label{chap:conclusion}

第~\ref{chap:overview}章与第~\ref{chap:sensitivity}章给出了统一 LP 框架与 $B^*$ 的灵敏度分析，
第~\ref{chap:experiments}章给出了实验设计。本章先明确方法有意划定的边界，再总结主要发现、局限与展望。

\section{方法边界}\label{sec:conc-boundary}

本文方法的能力与局限都源于同一组\textbf{约定}。超出约定的情形不在本方法讨论范围，列出如下：

\begin{itemize}
    \item \textbf{热的范围}：只讨论两个稳态分量——芯片峰值功耗与 PHY 额定动态功耗（第~\ref{sec:bg-thermal}节）。
    \textbf{翘曲、PDN 瞬态、瞬态热、良率、时钟}不在本方法讨论范围；
    \item \textbf{无阻塞的语义}：判定的是\textbf{无阻塞潜能}——最优自适应路由下物理资源足以支撑全部指定模式的同时交换，
    而非「保证」（第~\ref{sec:bg-nonblocking}节）。它是早期设计的早筛判据，假阴性不可接受、假阳性留给下游仿真消化；
    \item \textbf{对称性假设}：群论归约（附录 A）预设物理拓扑高度对称（vertex-transitive），
    对称性不足的拓扑不在讨论范围——这是以对称性换取线性规划可解性的取舍；
    \item \textbf{信号完整性}：经互连标准（UCIe/OIF-CEI 的 lane rate / reach / BER / 功耗档位）内嵌为约束，
    技术选型必须依赖已有互连标准，不单独建模；
    \item \textbf{die 缩放}：数值实验取 $\alpha_d=\beta_P=0$ 的退化特例（第~\ref{sec:ov-diescale}节），
    非零情形的灵敏度由第~\ref{chap:sensitivity}章的物理参数梯度覆盖。
\end{itemize}

这些边界是取舍而非缺陷：在边界之内，方法是严格的线性规划；
边界之外的情形需另起炉灶（缩小规模、换更强的无阻塞定义如 RNB，或引入专门仿真）。

\section{主要发现、局限与展望}\label{sec:conc-findings}

\textbf{主要发现。}对应第~\ref{sec:intro-contrib}节的三项贡献（群论归约是术、非核心贡献，不在其列）：

\begin{enumerate}
    \item \textbf{联立可行性判定}：把按物理位置组织的全部约束统一为同一组变量上的线性不等式组，
    将「多轮独立验证」替换为「单次联立判定」；精度升级（如热 L0/L1）只替换系数矩阵、不改变框架结构。
    \item \textbf{$B^*$ 统一筛选指标}：以 $B^* = f(\text{期待},\ \text{约束})$ 作为统一的设计点筛选指标，
    固定 $B$ 的可行性 LP 配合二分搜索，以对数次 LP 调用求得最大端口无阻塞带宽，
    使不同拓扑、不同技术选型在同一个标量上可比。
    \item \textbf{灵敏度分析}：在二分所得 $B^*$ 处用事后闭式 $\lambda_j = 1/(A_j\cdot\mathbf{L}^*)$ 提取绑定约束的影子价格，
    借助 Milgrom--Segal 包络定理给出 $B^*$ 对约束旋钮与期待旋钮的梯度，归一化后提供定量的投资优先级。
\end{enumerate}

三项贡献的共同落点是一句话：晶圆级交换机 DSE 中最困难的部分——多物理约束的耦合判断——
可以被一个线性规划干净地封装，并从「可行/不可行」提升到「改哪里、改多少」。

\textbf{局限。}
\begin{itemize}
    \item 实验数据尚在获取中，本章与第~\ref{chap:experiments}章所涉及的 $B^*$、Pareto 前沿、
    可扩展性与保真度的具体数值均待真实数据回填；
    \item die 缩放模型（$\alpha_d,\beta_P$）与组间 SerDes 约束的实现尚未接入当前 LP 求解路径；
    \item 布线约束依赖粗粒度的走线容量估计，完整布局布线数据有待与工艺标定对齐；
    \item 无阻塞潜能依赖拓扑对称性与「最优路由」假设，实际系统的实现效率需乘性松弛，并需与 cycle-accurate 仿真对标验证。
\end{itemize}

\textbf{展望。}
\begin{itemize}
    \item 接入 die 缩放与组间约束后，$B^*$ 将同时受 die 面积/峰值功耗随带宽的收缩与组间 C4 出口的双重约束，
    灵敏度分析可直接给出 $\alpha_d,\beta_P$ 的梯度指导；
    \item 把 RNB（可重排无阻塞，第~\ref{sec:bg-nonblocking}节）作为可行域下界纳入对比，
    为不依赖对称性假设的硬保证提供另一条路；
    \item 面向 $N>100$ 的规模，混合精度（L0 粗筛 + L1 精修）与分解求解是保持框架属性的可行路径；
    \item 以 cycle-accurate 仿真与工业设计数据对标，校准「潜能」到「保证」之间的实现效率乘性松弛。
\end{itemize}
